\let\negmedspace\undefined
\let\negthickspace\undefined
\documentclass[journal]{IEEEtran}
\usepackage[a5paper, margin=10mm, onecolumn]{geometry}
%\usepackage{lmodern} % Ensure lmodern is loaded for pdflatex
\usepackage{tfrupee} % Include tfrupee package

\setlength{\headheight}{1cm} % Set the height of the header box
\setlength{\headsep}{0mm}     % Set the distance between the header box and the top of the text

\usepackage{gvv-book}
\usepackage{gvv}
\usepackage{cite}
\usepackage{physics}
\usepackage{amsmath,amssymb,amsfonts,amsthm}
\usepackage{algorithmic}
\usepackage{graphicx}
\usepackage{textcomp}
\usepackage{xcolor}
\usepackage{txfonts}
\usepackage{listings}
\usepackage{enumitem}
\usepackage{mathtools}
\usepackage{gensymb}
\usepackage{comment}
\usepackage[breaklinks=true]{hyperref}
\usepackage{tkz-euclide} 
\usepackage{listings}
% \usepackage{gvv}                                        
\def\inputGnumericTable{}                                 
\usepackage[latin1]{inputenc}                                
\usepackage{color}                                            
\usepackage{array}                                            
\usepackage{longtable}                                       
\usepackage{calc}                                             
\usepackage{multirow}                                         
\usepackage{hhline}                                           
\usepackage{ifthen}                                           
\usepackage{lscape}
\begin{document}

\bibliographystyle{IEEEtran}

\title{5.2.25}
\author{EE25BTECH11022 - sankeerthan}
% \maketitle
% \newpage
% \bigskip
{\let\newpage\relax\maketitle}

\renewcommand{\thefigure}{\theenumi}
\renewcommand{\thetable}{\theenumi}
\setlength{\intextsep}{10pt} % Space between text and floats


\numberwithin{align}{enumi}
\numberwithin{figure}{enumi}
\renewcommand{\thetable}{\theenumi}

\textbf{Question}:Solve the following system of linear equations $ax+by =c$, $bx+ay =1+c$
\textbf{Solution:}  
Given  
\begin{align}
ax+by &= c \\
bx+ay &=1+c 
\end{align}
The matrix equation for a line is defined as
\begin{align}
    \vec{n}^\top\vec{x}=c
\end{align}
where $\vec{n}$ is the coefficient matrix and $\vec{x}=\myvec{x\\y}$
\begin{align}
       \myvec{a & b}\vec{x}&=c \\ 
       \myvec{b & a}\vec{x}=1+c
\end{align}
As a matrix equation
\begin{align}
\myvec{a&b\\b&a}\vec{x}&=\myvec{c \\ 1+c} \\
\myvec{a&b\\b&a}^\top\myvec{a&b\\b&a}\vec{x}&=\myvec{a&b\\b&a}^\top\myvec{c\\1+c} \\
\myvec{a&b\\b&a}\myvec{a&b\\b&a}\vec{x}&=\myvec{a&b\\b&a}\myvec{c\\1+c}   \\
\myvec{a^2+b^2 & 2ab \\2ab & a^2+b^2}\vec{x}&=\myvec{b+\brak{a+b}c \\a+\brak{a+b}c} \\
\vec{x}&=\frac{1}{(a^2-b^2)^2}\myvec{a^2+b^2 & -2ab \\-2ab& a^2+b^2}\myvec{b+\brak{a+b}c \\a+\brak{a+b}c} \\
\vec{x}&=\frac{1}{(a^2-b^2)^2} \myvec{\brak{a^2+b^2}\brak{b+\brak{a+b}c}-2ab\brak{a+\brak{a+b}c} \\-2ab\brak{b+\brak{a+b}c}+\brak{a^2+b^2}\brak{a+\brak{a+b}c}} \\
\vec{x}&=\frac{1}{(a^2-b^2)^2}\myvec{b\brak{a^2+b^2-2a^2}+ \brak{\brak{a+b}c}\brak{a^2+b^2-2ab} \\\brak{\brak{a+b}c}\brak{a^2+b^2-2ab}+ a\brak{a^2+b^2-2b^2}} \\
\myvec{x \\y}&=\frac{1}{(a^2-b^2)^2}\myvec{\brak{\brak{a+b}c}\brak{a-b}^2 -b\brak{a^2-b^2} \\\brak{\brak{a+b}c}\brak{a-b}^2 + a\brak{a^2-b^2}} \\
  \end{align}
Hence $x =\frac{c}{a+b}+\frac{b}{b^2-a^2},y=\frac{c}{a+b}+\frac{a}{a^2-b^2}$ 
is the solution for given system of linear equations

\begin{figure}
    \centering
    \includegraphics[width=1\columnwidth]{figs/points.png}
    \caption{}
    \label{}
\end{figure}
\end{document}